\begin{exbox}[线性失谐下的 Bloch 呼吸及零力极限]
前面的合成恒力项不必借助二态近似求解。忽略二次反冲，仅保留线性失谐和均匀最近邻耦合，令 $g>0$ 为角频率，
\begin{equation}
 \ii\dot c_\ell=\delta_B\ell c_\ell+g(c_{\ell-1}+c_{\ell+1}),
 \qquad c_\ell(0)=\delta_{\ell0}.
 \label{eq:feqo-bloch-linear-force-chain}
\end{equation}
求通道方差、复现周期和失谐趋零的极限。

\textbf{解。}在完整双向格 $\ell\in\ZZ$ 上，移位算符 $S|\ell\rangle=|\ell+1\rangle$ 满足 $SS^\dagger=S^\dagger S=I$。
令 $c_\ell=\ee^{-\ii\delta_B\ell t}a_\ell$，相互作用绘景的生成元为
$g(\ee^{\ii\delta_Bt}S+\ee^{-\ii\delta_Bt}S^\dagger)$。
因 $[S,S^\dagger]=0$，不同时间的生成元对易，时间排序精确消失：
\begin{equation*}
 U_I(t)=\exp[-\ii(A(t)S+A(t)^*S^\dagger)],\qquad
 A(t)=g\int_0^t\ee^{\ii\delta_Bu}\dd u
 =\frac{2g}{\delta_B}\ee^{\ii\delta_Bt/2}\sin(\delta_Bt/2).
\end{equation*}
Jacobi--Anger 展开给 $|c_\ell|^2=J_\ell^2(2|A|)$。
用已经证明的 Bessel 二阶矩，得到
\begin{equation}
 \langle\ell\rangle=0,\qquad
 \operatorname{Var}\ell=2|A|^2
 =8\left(\frac{g}{\delta_B}\right)^2\sin^2(\delta_Bt/2).
 \label{eq:feqo-bloch-breathing-variance}
\end{equation}
因此非零 $\delta_B$ 时分布以 $T_B=2\pi/|\delta_B|$ 完全复现，方差上限为 $8(g/\delta_B)^2$。
取极限 $\delta_B\to0$ 必须同时展开正弦，得到 $A\to gt$、$\operatorname{Var}\ell\to2g^2t^2$，没有发散。
这揭示恒力限制通道扩展的机制：各段跃迁振幅随失谐相位转动，在一个周期内抵消。
真实反冲的二次在位项会破坏此处的平移算符对易解；有限截断边界也会使 $[S,S^\dagger]\ne0$，必须先控制边界人口。

取 $r=g/|\delta_B|=0.75$，有界方差的峰值为 $8r^2=4.5$；在 $|\delta_B|t=4\pi$ 时恢复为零。
同一耦合而令力为零，经过相同时间后方差为 $2r^2(4\pi)^2=177.65$。
图中的两条曲线因此比较相同耦合和时长，不能把零力曲线当成另一种 Bloch 周期。
\end{exbox}

\begin{figure}[htbp]
 \centering
 \includegraphics[width=0.76\linewidth]{bloch_variance.pdf}
 \caption{线性失谐格的完整 Bessel 解。横轴为 $|\delta_B|t/\pi$，$r=g/|\delta_B|=0.75$。
 实线为 $8r^2\sin^2(\delta_Bt/2)$，虚线为在相同物理时长下的零力弹道扩展 $2(g t)^2$。
 二次反冲及边界截断未纳入这张图。}
 \label{fig:bloch-variance}
\end{figure}
